\documentclass[11pt]{article}
\usepackage[a4paper,margin=24mm]{geometry}
\usepackage{amsmath,amssymb,amsthm,booktabs,microtype}
\usepackage[colorlinks=true,urlcolor=blue,citecolor=blue]{hyperref}
\newtheorem{theorem}{Theorem}
\newtheorem{lemma}[theorem]{Lemma}
\newtheorem{corollary}[theorem]{Corollary}
\newcommand{\coef}[2]{[x^{#1}]#2}
\title{Unbounded consecutive failures of log-concavity\\from one fixed three-vertex branch}
\author{Research note prepared for Jan Mikul\'ik}
\date{27 September 2026}
\begin{document}
\maketitle
\begin{abstract}
Let $T_m$ have one root, $m$ children, and two pendant leaves at each child. We give an elementary proof that its independence polynomial has a consecutive interval of strict log-concavity failures of length at least
$2\log\log m/\log(3/2)-O(1)$.
The tree has diameter four, and every vertex except the root has degree at most three. Its repeated rooted branch is fixed. We also give an explicit finite sufficient criterion, requiring only integer comparisons. This is an independently derived structural specialization, not a claim of priority for unbounded consecutive failures, which were already established in work of Xiyao Yu.
\end{abstract}

\paragraph{Scope and prior work.}
Bautista-Ramos proved unbounded numbers of failures, without requiring consecutive indices~\cite{BR25}. Bautista-Ramos, Guill\'en-Galv\'an and G\'omez-Salgado constructed consecutive runs of specified small lengths~\cite{BR26}. Yu's September 2026 archive establishes unbounded runs and the substantially stronger unrestricted extremal lower bound $M(N)=\Omega(N^{1/3})$~\cite{Yu26}. The present note supplies a short proof for a fixed three-vertex module and diameter-four trees. Its $\log\log N$ bound does not improve Yu's unrestricted growth bound. Priority for this particular specialization has not been established. The argument below is self-contained; it does not resolve the tree unimodality conjecture.

\paragraph{Construction.}
For an integer $m\ge2$, the edges of $T_m$ are
\[
 \{vw_i,w_ia_i,w_ib_i:1\le i\le m\}.
\]
Thus $|V(T_m)|=3m+1$, its diameter is four, and $\alpha(T_m)=2m+1$.
Write $I_m(x)=\sum_k i_k(T_m)x^k$. Conditioning on the root gives
\begin{equation}\label{eq:I}
 I_m(x)=(1+3x+x^2)^m+x(1+x)^{2m}.
\end{equation}
The reflected polynomial is
\begin{equation}\label{eq:R}
 R_m(x)=x^{2m+1}I_m(1/x)=(1+x)^{2m}+x(1+3x+x^2)^m.
\end{equation}
Put $a_j=\coef{j}{R_m}=i_{2m+1-j}(T_m)$ and $r=3/2$.

\begin{theorem}[Finite criterion]\label{thm:finite}
If the positive integers $m,j$ satisfy
\begin{equation}\label{eq:criterion}
 j\ge48,\qquad
 \frac{72m}{j^2}\le r^j\le\frac m8,
 \qquad m\ge16(j+1)^3,
\end{equation}
then $a_j^2<a_{j-1}a_{j+1}$.
\end{theorem}

\begin{lemma}[Coefficient bounds]\label{lem:coeff}
Define
\[
 b_0=1,\qquad b_j=\frac{(2m)^j}{j!}
       +\frac{(3m)^{j-1}}{(j-1)!}\quad(j\ge1).
\]
For $0\le j\le m$,
\begin{equation}\label{eq:bounds}
 \left(1-\frac{j(j-1)}{2m}\right)b_j\le a_j\le b_j.
\end{equation}
\end{lemma}
\begin{proof}
For $0\le k\le m$,
\[
 \binom mk3^k\le\coef{k}{(1+3x+x^2)^m}
                 \le\frac{(3m)^k}{k!}.
\]
The lower bound uses the terms choosing $3x$ from $k$ distinct factors. For the upper bound factor $1+3x+x^2=(1+\lambda_+x)(1+\lambda_-x)$, where $\lambda_\pm=(3\pm\sqrt5)/2>0$. The $k$th elementary symmetric polynomial in nonnegative weights is at most the $k$th power of their sum divided by $k!$. Also
\[
 \binom mk=\frac{m^k}{k!}\prod_{s=0}^{k-1}(1-s/m)
 \ge\frac{m^k}{k!}\left(1-\frac{k(k-1)}{2m}\right).
\]
Apply these bounds with $k=j-1$, and the analogous bounds to
$\binom{2m}{j}$. Their sum gives~\eqref{eq:bounds}; $j=0$ is immediate.
\end{proof}

\begin{proof}[Proof of Theorem~\ref{thm:finite}]
For $j\ge1$ put $t=j r^j/(3m)$, so $b_j=(2m)^j(1+t)/j!$.
The adjacent values of this parameter are $t(j-1)/(jr)$ and
$t(j+1)r/j$. Direct algebra gives the exact identity
\begin{equation}\label{eq:curvature}
 \frac{b_{j-1}b_{j+1}}{b_j^2}-1
 =\frac{t}{(j+1)(1+t)^2}
 \left(\frac{j-7}{6}-\frac1t-\left(1+\frac1j\right)t\right).
\end{equation}
The middle conditions in~\eqref{eq:criterion} imply
$24/j\le t\le j/24$. Hence, for $j\ge48$,
\[
 \frac{j-7}{6}-\frac1t-\left(1+\frac1j\right)t
 \ge\frac{j}{12}-\frac{29}{24}\ge\frac j{24}.
\]
On that same interval,
\[
 \frac{t}{(1+t)^2}\ge
 \frac{24/j}{(1+24/j)^2}\ge\frac{32}{3j}.
\]
Therefore~\eqref{eq:curvature} is at least $4/(9(j+1))\ge1/(4j)$.
Let $\delta=(j+1)^2/(2m)$. Lemma~\ref{lem:coeff}, applied at the three
indices, now yields
\[
 \frac{a_{j-1}a_{j+1}}{a_j^2}
 \ge(1-\delta)^2\left(1+\frac1{4j}\right)
 \ge\left(1-\frac1{16(j+1)}\right)
       \left(1+\frac1{4j}\right)>1.
\]
Here $m\ge16(j+1)^3$ ensures both the applicability of the lemma and
$2\delta\le1/(16(j+1))$. This proves the strict inequality.
\end{proof}

\begin{corollary}[Unbounded runs with a fixed branch]\label{cor:run}
For all sufficiently large $m$, the polynomial $I_m$ has a consecutive
interval of strict log-concavity failures of length
\[
 \frac{2}{\log(3/2)}\log\log m+O(1).
\]
In particular the longest such interval is at least this quantity. The
coefficient offsets from degree $2m+1$ lie between
\[
 \frac{\log m-2\log\log m+O(1)}{\log(3/2)}
 \quad\text{and}\quad
 \frac{\log m+O(1)}{\log(3/2)}.
\]
\end{corollary}
\begin{proof}
Let $\lambda=\log(3/2)$, and let $u_m>0$ be the unique real solution of
$u_m^2e^{\lambda u_m}=72m$. Put $v_m=\log(m/8)/\lambda$.
Taking logarithms and substituting back gives
\[
 u_m=\frac{\log m-2\log\log m+\log72+2\log\lambda+o(1)}{\lambda},
 \qquad v_m=\frac{\log m-\log8}{\lambda}.
\]
Every integer $j\in[\lceil u_m\rceil,\lfloor v_m\rfloor]$ satisfies the
middle inequalities of~\eqref{eq:criterion}, since $j^2r^j$ is increasing.
For sufficiently large $m$ the interval is nonempty, its lower endpoint
is at least $48$, and $m\ge16(\lfloor v_m\rfloor+1)^3$.
Theorem~\ref{thm:finite} applies throughout. The number of integers in the
interval equals $(2/\lambda)\log\log m+O(1)$. Reflection reverses the
interval but preserves the strict log-concavity inequalities.
\end{proof}

\paragraph{Why the effect persists.}
The two terms in~\eqref{eq:R} have coefficients approximately
$(2m)^j/j!$ and $(3m)^{j-1}/(j-1)!$. Their ratio is
$j(3/2)^j/(3m)$. The crossover moves to offsets of order $\log m$,
where the log-concavity of the factorial denominator has size $1/j$.
The positive curvature from the mixture exceeds that size over an interval
of width proportional to $\log j$. Formula~\eqref{eq:curvature} and the
finite bounds make this explanation rigorous.

\paragraph{Exact checks.}
The accompanying standard-library Python program calculates coefficient
prefixes with arbitrary-precision integers. The displayed intervals were
checked directly using $a_{j-1}a_{j+1}-a_j^2>0$; no floating-point sign
checks are used. The intervals are maximal within the computed prefix,
not a claim to enumerate all failures in the full polynomial.
\begin{center}
\begin{tabular}{@{}rrr@{}}
\toprule
$m$ & Reflected failure offsets & Consecutive length\\
\midrule
$2^{16}$ & $22$--$24$ & $3$\\
$2^{32}$ & $44$--$52$ & $9$\\
$2^{64}$ & $95$--$107$ & $13$\\
$2^{128}$ & $200$--$217$ & $18$\\
$2^{256}$ & $416$--$435$ & $20$\\
\bottomrule
\end{tabular}
\end{center}
Writing $p_k=\coef{k}{(1+3x+x^2)^m}$ gives the exact recurrence
\[
 p_0=1,\quad p_1=3m,\qquad
 (k+1)p_{k+1}=3(m-k)p_k+(2m-k+1)p_{k-1}.
\]
Then $a_0=1$ and $a_j=\binom{2m}{j}+p_{j-1}$. The verifier also compares
this recurrence against direct polynomial multiplication on small trees,
and checks the finite criterion independently of the observed intervals.

\paragraph{Status.}
The mathematical claim is exactly Theorem~\ref{thm:finite} and its
corollary, with the proof printed above. The numerical checks are
supplementary. This draft has not undergone independent expert review.
The unrestricted existence theorem belongs to the prior literature;
novelty of the fixed-module, diameter-four specialization remains to be
checked beyond the sources consulted. The failures occur in the far tail
and do not imply non-unimodality.

\begin{thebibliography}{9}\small
\bibitem{BR25} C. Bautista-Ramos,
\emph{Multiple breaks of log-concavity in the independence polynomials of trees},
arXiv:2511.00334 (2025).
\url{https://arxiv.org/abs/2511.00334}.
\bibitem{BR26} C. Bautista-Ramos, C. Guill\'en-Galv\'an and P. G\'omez-Salgado,
\emph{Linear recurrences for non-log-concave independence polynomials of trees},
Graphs and Combinatorics \textbf{42}, 59 (2026).
\url{https://doi.org/10.1007/s00373-026-03054-4}.
\bibitem{Yu26} X. Yu,
\emph{Unbounded runs of non-log-concavity in independence polynomials of trees},
manuscript in reproducibility archive, version 1.1.0, 3 September 2026.
\url{https://doi.org/10.5281/zenodo.22263331}.
\end{thebibliography}
\end{document}
